% Generated by latex-iot-paper-generator | Dataset(s): CIC-IoT, WUSTL-IIoT, TON-IoT, Bot-IoT, UNSW-NB15 | Date: 2026-04-29
% Kerckhoffs robustness data from: experiments/results/kerckhoffs/kerckhoffs-results-2026-04-28-14-52-13.json
% Zero-day anon ablation data: 1-cic-iot/results/anon/ and 2-wustl-iiot/results/anon/

\appendix

%=======================================================================
\section{Detailed Per-Class Metrics}\label{app:detailed}
%=======================================================================

%-----------------------------------------------------------------------
\subsection{CIC-IoT2023 — LLM Multiclass Full Report}
%-----------------------------------------------------------------------

Table~\ref{tab:app_cic_llm_full} reproduces the full classification report
extracted from \texttt{result-llm-multiclass-100000.txt}. The confusion matrix
is displayed in Figure~\ref{fig:cm_cic_llm}.

\begin{table*}[htbp]
\centering
\caption{Full LLM multiclass classification report, CIC-IoT2023
  (100{,}000-sample corpus).}
\label{tab:app_cic_llm_full}
\begin{tabular}{lcccc}
\toprule
\textbf{Class} & \textbf{Precision} & \textbf{Recall} & \textbf{F1-Score} & \textbf{Support} \\
\midrule
BenignTraffic      & 0.2865 & 0.3426 & 0.3120 & 470 \\
DDoS-Flood         & 0.5163 & 0.4383 & 0.4741 & 470 \\
DDoS-Fragmentation & 0.7482 & 0.9400 & 0.8332 & 433 \\
DNS\_Spoofing      & 0.1065 & 0.2222 & 0.1440 &  81 \\
DoS-Flood          & 0.6370 & 0.1979 & 0.3019 & 470 \\
MITM-ArpSpoofing   & 0.3540 & 0.3053 & 0.3279 & 131 \\
Mirai              & 0.8068 & 0.9064 & 0.8537 & 470 \\
Recon              & 0.3024 & 0.4397 & 0.3584 & 141 \\
\midrule
Accuracy           & \multicolumn{4}{c}{0.5300} \\
Macro Avg          & 0.4697 & 0.4740 & 0.4506 & 2666 \\
Weighted Avg       & 0.5542 & 0.5300 & 0.5171 & 2666 \\
\bottomrule
\end{tabular}
\end{table*}

%-----------------------------------------------------------------------
\subsection{CIC-IoT2023 — Random Forest Multiclass Full Report}
%-----------------------------------------------------------------------

\begin{table*}[htbp]
\centering
\caption{Full Random Forest multiclass classification report, CIC-IoT2023.}
\label{tab:app_cic_rf_full}
\begin{tabular}{lcccc}
\toprule
\textbf{Class} & \textbf{Precision} & \textbf{Recall} & \textbf{F1-Score} & \textbf{Support} \\
\midrule
BenignTraffic      & 0.8577 & 0.9617 & 0.9067 & 470 \\
DDoS-Flood         & 1.0000 & 0.9957 & 0.9979 & 470 \\
DDoS-Fragmentation & 0.9931 & 0.9977 & 0.9954 & 433 \\
DNS\_Spoofing      & 0.8367 & 0.5062 & 0.6308 &  81 \\
DoS-Flood          & 0.9979 & 1.0000 & 0.9989 & 470 \\
MITM-ArpSpoofing   & 0.8684 & 0.7557 & 0.8082 & 131 \\
Mirai              & 1.0000 & 0.9957 & 0.9979 & 470 \\
Recon              & 0.7836 & 0.7447 & 0.7636 & 141 \\
\midrule
Accuracy           & \multicolumn{4}{c}{0.9509} \\
Macro Avg          & 0.9172 & 0.8697 & 0.8874 & 2666 \\
Weighted Avg       & 0.9505 & 0.9509 & 0.9487 & 2666 \\
\bottomrule
\end{tabular}
\end{table*}

%-----------------------------------------------------------------------
\subsection{WUSTL-IIoT — Both Zero-Day Runs Compared}\label{app:wustl_zeroday}
%-----------------------------------------------------------------------

Two independent zero-day runs were conducted for WUSTL-IIoT with
\texttt{Reconn} as the withheld class (seed~42). Table~\ref{tab:zeroday_wustl_both}
compares both runs. The second run (\texttt{2026-04-20-17-58-16}) uses a
different rule set based on rate features and achieves ZDR~$= 1.0$; the first
run (\texttt{2026-04-20-17-07-34}) uses packet-count features and also reaches
ZDR~$= 1.0$ (reported independently). The main paper uses the higher-fidelity
run sourced from \texttt{zeroday-Reconn-seed42-2026-04-20-17-58-16.json}.

\begin{table*}[htbp]
\centering
\caption{WUSTL-IIoT zero-day runs comparison (withheld: Reconn, seed 42,
  8{,}240 zero-day samples).}
\label{tab:zeroday_wustl_both}
\begin{tabular}{llccc}
\toprule
\textbf{Run Timestamp} & \textbf{Top Rules (features)} & \textbf{Known-F1} & \textbf{ZDR} \\
\midrule
2026-04-20-17-07-34 & SrcRate, DstRate, Rate, Dur, TcpRtt & 0.7854 & 1.0000 \\
2026-04-20-17-58-16 & SrcPkts, DstPkts, TotPkts, DstBytes, SrcBytes & 0.8005 & 0.9989 \\
\bottomrule
\end{tabular}
\end{table*}

%-----------------------------------------------------------------------
\subsection{Anonymization Ablation --- Per-Seed Breakdown}
%-----------------------------------------------------------------------

Tables~\ref{tab:anon_perseed_cic}--\ref{tab:anon_perseed_botiot} report
per-seed F1 values for the anonymization ablation to allow verification of
variance. The mean Jaccard across seeds measures how consistently each
condition selects the same $k=5$ features.

\begin{table*}[htbp]
\centering
\caption{CIC-IoT2023 anonymization ablation --- per-seed F1
  (binary attack/normal, $n{=}100{,}000$).}
\label{tab:anon_perseed_cic}
\begin{tabular}{lccc|ccc}
\toprule
 & \multicolumn{3}{c|}{\textbf{F1 (Original)}} & \multicolumn{3}{c}{\textbf{F1 (Anonymized)}} \\
\textbf{Seed} & \textbf{Seed 42} & \textbf{Seed 123} & \textbf{Seed 456} &
               \textbf{Seed 42} & \textbf{Seed 123} & \textbf{Seed 456} \\
\midrule
Value & 0.6642 & 0.8532 & 0.5955 & 0.8585 & 0.6656 & 0.6538 \\
\midrule
\multicolumn{7}{l}{Mean Jaccard (orig vs.\ anon feature sets): 0.889 $\pm$ 0.157}\\
\multicolumn{7}{l}{Per-seed Jaccard: [1.000, 1.000, 0.667]}\\
\bottomrule
\end{tabular}
\end{table*}

\begin{table*}[htbp]
\centering
\caption{WUSTL-IIoT anonymization ablation --- per-seed F1
  (binary attack/normal).}
\label{tab:anon_perseed_wustl}
\begin{tabular}{lccc|ccc}
\toprule
 & \multicolumn{3}{c|}{\textbf{F1 (Original)}} & \multicolumn{3}{c}{\textbf{F1 (Anonymized)}} \\
\textbf{Seed} & \textbf{Seed 42} & \textbf{Seed 123} & \textbf{Seed 456} &
               \textbf{Seed 42} & \textbf{Seed 123} & \textbf{Seed 456} \\
\midrule
Value & 0.9530 & 0.8134 & 0.7709 & 0.6062 & 0.7318 & 0.6653 \\
\midrule
\multicolumn{7}{l}{Mean Jaccard (orig vs.\ anon feature sets): 0.083 $\pm$ 0.118}\\
\multicolumn{7}{l}{Per-seed Jaccard: [0.000, 0.250, 0.000]}\\
\bottomrule
\end{tabular}
\end{table*}

\begin{table*}[htbp]
\centering
\caption{TON\_IoT anonymization ablation --- per-seed F1 (binary attack/normal).}
\label{tab:anon_perseed_toniot}
\begin{tabular}{lccc|ccc}
\toprule
 & \multicolumn{3}{c|}{\textbf{F1 (Original)}} & \multicolumn{3}{c}{\textbf{F1 (Anonymized)}} \\
\textbf{Seed} & \textbf{Seed 42} & \textbf{Seed 123} & \textbf{Seed 456} &
               \textbf{Seed 42} & \textbf{Seed 123} & \textbf{Seed 456} \\
\midrule
Value & 0.3414 & 0.7167 & 0.6727 & 0.6107 & 0.7121 & 0.7001 \\
\midrule
\multicolumn{7}{l}{Mean Jaccard (orig vs.\ anon feature sets): 0.310 $\pm$ 0.084}\\
\multicolumn{7}{l}{Per-seed Jaccard: [0.250, 0.250, 0.429]}\\
\bottomrule
\end{tabular}
\end{table*}

\begin{table*}[htbp]
\centering
\caption{Bot-IoT anonymization ablation --- per-seed F1 (binary attack/normal).}
\label{tab:anon_perseed_botiot}
\begin{tabular}{lccc|ccc}
\toprule
 & \multicolumn{3}{c|}{\textbf{F1 (Original)}} & \multicolumn{3}{c}{\textbf{F1 (Anonymized)}} \\
\textbf{Seed} & \textbf{Seed 42} & \textbf{Seed 123} & \textbf{Seed 456} &
               \textbf{Seed 42} & \textbf{Seed 123} & \textbf{Seed 456} \\
\midrule
Value & 0.9120 & 0.8509 & 0.8578 & 0.8845 & 0.9052 & 0.8912 \\
\midrule
\multicolumn{7}{l}{Mean Jaccard (orig vs.\ anon feature sets): 0.587 $\pm$ 0.112}\\
\multicolumn{7}{l}{Per-seed Jaccard: [0.429, 0.667, 0.667]}\\
\bottomrule
\end{tabular}
\end{table*}

%-----------------------------------------------------------------------
\subsection{Zero-Day Anonymization Ablation}\label{app:zeroday_anon}
%-----------------------------------------------------------------------

Table~\ref{tab:zeroday_anon_comparison} reports the impact of anonymization
on zero-day detection capability for the two datasets where a dedicated
zero-day anonymization ablation file was produced
(\texttt{zeroday-Mirai-udpplain-anon-ablation} for CIC-IoT2023 and
\texttt{zeroday-Reconn-anon-ablation} for WUSTL-IIoT).

\begin{table*}[htbp]
\centering
\caption{Zero-day detection under original versus anonymized feature names.
  FPR~=~False Positive Rate on benign traffic;
  ZDR~=~Zero-Day Detection Rate on withheld class;
  Known-F1~=~macro F1 on remaining known classes.
  $\Delta$ZDR~=~ZDR$_{\text{anon}}$~$-$~ZDR$_{\text{orig}}$.}
\label{tab:zeroday_anon_comparison}
\begin{tabular}{lllccc|ccc|c}
\toprule
\textbf{Dataset} & \textbf{Withheld} & \textbf{Condition} &
  \textbf{Known-F1} & \textbf{ZDR} & \textbf{FPR} &
  \multicolumn{3}{c|}{\textbf{Top Rules (features)}} &
  \textbf{$\Delta$ZDR} \\
\midrule
CIC-IoT2023  & Mirai-udpplain &
  Original   & [INSERT: known\_f1\_orig\_cic\_zd\_anon] & [INSERT: zdr\_orig] & [INSERT: fpr\_orig] &
  \multicolumn{3}{c|}{flow\_duration, Header\_Length, Rate, Tot sum, Covariance} & \\
             &  &
  Anonymized & 0.9531 & [INSERT: zdr\_anon\_ciciot] & [INSERT: fpr\_anon\_ciciot] &
  \multicolumn{3}{c|}{f0 (flow\_duration), f1 (Header\_Length), f4 (Rate), f30 (Tot sum), f40 (Covariance)} &
  [INSERT: delta\_zdr\_ciciot] \\
\midrule
WUSTL-IIoT   & Reconn &
  Original   & 0.8560 & 0.8140 & 0.0165 &
  \multicolumn{3}{c|}{SrcPkts$<$5, DstPkts$\leq$1, TotBytes$<$500, SrcLoad$>$300k, Rate$>$20k} & \\
             &  &
  Anonymized & 0.8631 & \textbf{0.9717} & 0.0278 &
  \multicolumn{3}{c|}{Sport$<$50k, Dport$<$502, SrcPkts$<$10, DstPkts$<$8, SrcLoad$>$100k} &
  $+$0.1577 \\
\bottomrule
\end{tabular}
\end{table*}

A notable finding is that for WUSTL-IIoT, the anonymized condition achieves
\emph{higher} ZDR ($0.972$) than the original ($0.814$) at a slightly higher
FPR ($0.028$ vs.\ $0.017$). This counterintuitive result suggests that without
semantic cues, the LLM gravitates toward port-range rules that coincide more
precisely with the reconnaissance scanning pattern of the withheld class.

%=======================================================================
\section{Confusion Matrices}\label{app:confusion}
%=======================================================================

[INSERT: confusion\_matrix\_CICIoT2023\_LLM]

\noindent\textit{Note:} The raw confusion matrix values for the CIC-IoT2023
LLM multiclass run are as follows (rows: true class, columns: predicted class,
order: BenignTraffic, DDoS-Flood, DDoS-Fragmentation, DNS\_Spoofing,
DoS-Flood, MITM-ArpSpoofing, Mirai, Recon):

\vspace{0.5em}
\noindent\texttt{
[[161,\;4,\;104,\;105,\;0,\;32,\;1,\;63],} \\
\texttt{\phantom{[}[177,\;206,\;0,\;0,\;52,\;0,\;26,\;9],} \\
\texttt{\phantom{[}[10,\;0,\;407,\;1,\;0,\;4,\;1,\;10],} \\
\texttt{\phantom{[}[26,\;0,\;7,\;18,\;0,\;14,\;1,\;15],} \\
\texttt{\phantom{[}[116,\;173,\;0,\;0,\;93,\;4,\;72,\;12],} \\
\texttt{\phantom{[}[38,\;2,\;18,\;15,\;0,\;40,\;1,\;17],} \\
\texttt{\phantom{[}[10,\;0,\;0,\;1,\;0,\;16,\;426,\;17],} \\
\texttt{\phantom{[}[24,\;14,\;8,\;29,\;1,\;3,\;0,\;62]]}
\vspace{0.5em}

\noindent The RF multiclass confusion matrix (same class order):

\vspace{0.5em}
\noindent\texttt{
[[452,\;0,\;0,\;3,\;0,\;1,\;0,\;14],} \\
\texttt{\phantom{[}[0,\;468,\;1,\;0,\;1,\;0,\;0,\;0],} \\
\texttt{\phantom{[}[1,\;0,\;432,\;0,\;0,\;0,\;0,\;0],} \\
\texttt{\phantom{[}[19,\;0,\;0,\;41,\;0,\;11,\;0,\;10],} \\
\texttt{\phantom{[}[0,\;0,\;0,\;0,\;470,\;0,\;0,\;0],} \\
\texttt{\phantom{[}[24,\;0,\;0,\;3,\;0,\;99,\;0,\;5],} \\
\texttt{\phantom{[}[0,\;0,\;2,\;0,\;0,\;0,\;468,\;0],} \\
\texttt{\phantom{[}[31,\;0,\;0,\;2,\;0,\;3,\;0,\;105]]}
\vspace{0.5em}

[INSERT: confusion\_matrix\_WUSTL\_IIoT]

[INSERT: confusion\_matrix\_TON\_IoT]

[INSERT: confusion\_matrix\_Bot\_IoT]

[INSERT: confusion\_matrix\_UNSW\_NB15]

%=======================================================================
\section{Threshold Sensitivity Analysis}\label{app:threshold_full}
%=======================================================================

All threshold sensitivity figures are reproduced here for completeness.
Figures~\ref{fig:app_threshold_mitm} and~\ref{fig:app_threshold_mirai} cover
the two CIC-IoT2023 zero-day scenarios. Figures~\ref{fig:app_threshold_wustl},
\ref{fig:app_threshold_toniot}, and~\ref{fig:app_threshold_botiot} cover the
remaining datasets. Cross-references to the main text are provided in
Table~\ref{tab:threshold_xref}.

\begin{table}[htbp]
\centering
\caption{Cross-reference between threshold figures and main text.}
\label{tab:threshold_xref}
\begin{tabular}{lll}
\toprule
\textbf{Dataset} & \textbf{Withheld Class} & \textbf{Figure} \\
\midrule
CIC-IoT2023  & MITM-ArpSpoofing & Fig.~\ref{fig:app_threshold_mitm} \\
CIC-IoT2023  & Mirai-udpplain   & Fig.~\ref{fig:app_threshold_mirai} \\
WUSTL-IIoT   & Reconn           & Fig.~\ref{fig:app_threshold_wustl} \\
TON\_IoT     & ransomware       & Fig.~\ref{fig:app_threshold_toniot} \\
Bot-IoT      & Service\_Scan    & Fig.~\ref{fig:app_threshold_botiot} \\
\bottomrule
\end{tabular}
\end{table}

\begin{figure}[htbp]
  \centering
  \includegraphics[width=\columnwidth]{../iot-llm-hids-pg/1-cic-iot/results/zeroday-MITM-ArpSpoofing-thresholds.png}
  \caption{Full threshold sensitivity curve: CIC-IoT2023, MITM-ArpSpoofing
    withheld (11{,}211 samples, seed 42).
    Known-F1 on remaining classes: 0.8648. ZDR at best threshold: 0.0218.}
  \label{fig:app_threshold_mitm}
\end{figure}

\begin{figure}[htbp]
  \centering
  \includegraphics[width=\columnwidth]{../iot-llm-hids-pg/1-cic-iot/results/zeroday-Mirai-udpplain-thresholds.png}
  \caption{Full threshold sensitivity curve: CIC-IoT2023, Mirai-udpplain
    withheld (31{,}864 samples, seed 42).
    Known-F1 on remaining classes: 0.8301. ZDR at best threshold: 0.0004.}
  \label{fig:app_threshold_mirai}
\end{figure}

\begin{figure}[htbp]
  \centering
  \includegraphics[width=\columnwidth]{../iot-llm-hids-pg/2-wustl-iiot/results/zeroday-Reconn-thresholds.png}
  \caption{Full threshold sensitivity curve: WUSTL-IIoT, Reconn withheld
    (8{,}240 samples, seed 42).
    Known-F1 on remaining classes: 0.8005. ZDR at best threshold: 0.9989.}
  \label{fig:app_threshold_wustl}
\end{figure}

\begin{figure}[htbp]
  \centering
  \includegraphics[width=\columnwidth]{../iot-llm-hids-pg/3-ton-iot/results/zeroday-ransomware-thresholds.png}
  \caption{Full threshold sensitivity curve: TON\_IoT, ransomware withheld
    (14{,}735 samples, seed 42).
    Known-F1 on remaining classes: 0.6745. ZDR at best threshold: 0.9999.}
  \label{fig:app_threshold_toniot}
\end{figure}

\begin{figure}[htbp]
  \centering
  \includegraphics[width=\columnwidth]{../iot-llm-hids-pg/4-bot-iot/results/zeroday-Service_Scan-thresholds.png}
  \caption{Full threshold sensitivity curve: Bot-IoT, Service\_Scan withheld
    (58{,}626 samples, seed 42).
    Known-F1 on remaining classes: 0.9662. ZDR at best threshold: 0.6490.}
  \label{fig:app_threshold_botiot}
\end{figure}

%=======================================================================
\section{Robustness and Adversarial Hardness Analysis}\label{app:robustness}
%=======================================================================

Table~\ref{tab:robustness} reports the critical perturbation budget
$\varepsilon_{0.5}$, defined as the minimum adversarial perturbation (as a
fraction of feature IQR) required to evade 50\% of detections, for all five
detector conditions across two datasets (CIC-IoT2023 and WUSTL-IIoT).
Values are extracted from the most recent Kerckhoffs experiment run
(\texttt{kerckhoffs-results-2026-04-28-14-52-13.json}, seeds~42, 123, 456).
A \emph{higher} $\varepsilon_{0.5}$ indicates a \emph{harder-to-evade} detector.
The figure \texttt{robustness\_curves.png} plots the full ESR curves.

\begin{table*}[htbp]
\centering
\caption{Adversarial hardness: critical perturbation budget $\varepsilon_{0.5}$
  (mean $\pm$ std, 3 seeds) at which the Exploitability Sweep Rate reaches 0.5.
  Higher is harder to evade. LLM\_orig: LLM with semantic feature names;
  LLM\_anon: LLM with anonymized names; DT\_d3/d5: Decision Tree depth 3/5;
  RF\_d5: Random Forest depth 5.}
\label{tab:robustness}
\begin{tabular}{llccccc}
\toprule
\textbf{Dataset} & \textbf{Withheld} &
  \textbf{LLM\_orig} & \textbf{LLM\_anon} & \textbf{DT\_d3} &
  \textbf{DT\_d5} & \textbf{RF\_d5} \\
\midrule
CIC-IoT2023 & Mirai-udpplain &
  $0.00781 \pm 0.000046$ &
  $\mathbf{0.00213 \pm 0.000014}$ &
  $0.00290 \pm 0.000112$ &
  $0.00290 \pm 0.000112$ &
  $0.01024 \pm 0.003471$ \\
WUSTL-IIoT  & Reconn &
  $\mathbf{0.00180 \pm 0.000017}$ &
  $0.01789 \pm 0.000112$ &
  $0.03494 \pm 0.0000021$ &
  $0.00050 \pm 0.000000030$ &
  $0.00051 \pm 0.000016$ \\
\bottomrule
\end{tabular}
\end{table*}

\textbf{Key observations.}

\textit{CIC-IoT2023.} The RF\_d5 detector shows the highest mean
$\varepsilon_{0.5}$ ($0.010$) but also the highest variance across seeds, suggesting
inconsistent robustness. LLM\_orig and RF\_d5 are comparably hard to evade;
LLM\_anon is the \emph{easiest} to evade on this dataset ($\varepsilon_{0.5} =
0.00213$), requiring only 0.2\% IQR perturbation to evade half of detections.

\textit{WUSTL-IIoT.} The pattern reverses: LLM\_orig is the \emph{hardest}
to evade ($\varepsilon_{0.5} = 0.00180$), while DT\_d3 requires the largest budget
to break ($\varepsilon_{0.5} = 0.035$). DT\_d5 and RF\_d5 are the most fragile,
collapsing at $\varepsilon_{0.5} \approx 0.0005$. LLM\_anon occupies an
intermediate position ($\varepsilon_{0.5} = 0.018$), consistent with the
observation that anonymized rules use port-range conditions that are harder to
satisfy simultaneously under a small perturbation budget.

\begin{figure}[htbp]
  \centering
  \includegraphics[width=\columnwidth]{../iot-llm-hids-pg/experiments/results/kerckhoffs/robustness_curves.png}
  \caption{Full ESR curves for all five detector conditions on CIC-IoT2023
    (Mirai-udpplain withheld) and WUSTL-IIoT (Reconn withheld).
    Each curve shows the fraction of flagged samples that evade detection
    at increasing perturbation budgets $\varepsilon$.}
  \label{fig:robustness_curves}
\end{figure}

\begin{figure}[htbp]
  \centering
  \includegraphics[width=\columnwidth]{../iot-llm-hids-pg/experiments/results/kerckhoffs/feature_targeting.png}
  \caption{Feature targeting heatmap: number of flagged samples whose
    evasion path modifies each feature, broken down by detector condition.
    Dominance of single features (e.g.\ \texttt{rst\_count} for DT\_d3,
    \texttt{TotBytes} for LLM\_orig on WUSTL-IIoT) reveals brittle
    single-feature rules; LLM-generated rules typically distribute
    evasion load across multiple features.}
  \label{fig:feature_targeting}
\end{figure}
